\section{Related Work}
\label{sec:related}

\para{Hidden state capacity and memory tokens.}
\citet{kuratov2025cramming} is the most directly related work, demonstrating that a single trainable \memvec vector can encode up to 1,568 tokens when decoded through a frozen Llama-3.1-8B model. They establish a theoretical upper bound of $\dmodel \times b / \ceil{\log_2 \vocabsize}$ tokens and show that empirical capacity utilization is only 5--30\% of this maximum. Their approach optimizes a single vector via gradient descent with a frozen LLM decoder. We extend their analysis by systematically comparing this full-model decoding against four simpler encoding schemes, revealing how much of the capacity comes from the encoding itself versus the decoder.

\para{Superposition and compressed representations.}
\citet{elhage2022superposition} showed that neural networks store more features than they have dimensions by exploiting feature sparsity, with features organizing as geometric polytopes in weight space. This superposition phenomenon enables networks to represent $m > d$ features in $d$ dimensions when features are sparse. We test whether superposition extends to token ID encoding---where signals are dense and high-entropy rather than sparse and binary---and find that it fails entirely, clarifying the boundary conditions for superposition.

\para{Decoding hidden states.}
The logit lens \citep{dar2023analyzing} projects intermediate hidden states through the unembedding matrix to obtain vocabulary distributions, revealing what information is accessible at each layer. \citet{belrose2023tuned} improved this with the tuned lens, training per-layer affine transformations that account for representational drift between layers. \citet{pal2023future} extended these ideas with Future Lens, showing that a single hidden state encodes information about tokens at positions $t+2$, $t+3$, and beyond. Patchscopes \citep{ghandeharioun2024patchscopes} unifies these approaches under a framework that patches hidden states between inference passes. Our unembedding-based encoding scheme (\secref{sec:method_unembed}) builds on these insights by using the frozen unembedding matrix as a fixed decoder and learning only a linear projection from subspace chunks.

\para{Information content of hidden states.}
\citet{geva2021transformer} showed that transformer feed-forward layers operate as key-value memories, with lower layers storing shallow patterns and upper layers storing semantic information. \citet{voita2020information} proposed information-theoretic probing based on minimum description length, providing a principled framework for measuring how much information representations encode about specific properties. \citet{park2023linear} formalized the linear representation hypothesis---that high-level concepts correspond to linear directions in representation space---while \citet{zou2023representation} demonstrated that these directions can be read and controlled for safety properties. These findings inform our geometric encoding schemes, which exploit the directional structure of the embedding space.

\para{Capacity bounds and the softmax bottleneck.}
\citet{yang2018breaking} identified the softmax bottleneck: the unembedding matrix has rank $\dmodel \ll \vocabsize$, limiting the expressiveness of hidden-state-to-vocabulary mappings. \citet{kajitsuka2024optimal} proved that transformers can memorize $N$ input sequences with $O(\sqrt{N})$ parameters, establishing theoretical capacity bounds. \citet{morris2025memorize} showed that GPT-style transformers store approximately 3.5--4 bits per model parameter. These theoretical results contextualize our empirical measurements and help explain why model-based decoders achieve far less than the raw bit capacity of the hidden state.

\para{Hyperspherical representations.}
\citet{loshchilov2024ngpt} proposed nGPT, which constrains all representations to unit norm on a hypersphere, making direction the primary information carrier. This is directly relevant to our vector walk encoding scheme, where magnitude encodes token identity and direction encodes position. The success of nGPT suggests that angular geometry is a natural and efficient structure for encoding information in hidden states.
